\section{Introducción}
La abundancia bariónica actual conserva información sobre el exceso de materia respecto de antimateria, pero reconstruir las cantidades que coexistieron durante la aniquilación exige establecer qué información adicional aporta un modelo térmico. Dos inventarios con igual diferencia entre bariones y antibariones pueden contener cantidades de pares muy distintas. Por ello, la pregunta retrospectiva debe distinguir el residuo observado de la población total inferida a una temperatura especificada.

La literatura ofrece referencias concretas para situar esa comparación. \citet{Satarov2013} describen abundancias de bariones y antibariones casi iguales por encima de aproximadamente 50 MeV y una disminución rápida de la proporción de antibariones al enfriarse el plasma por debajo de esa escala. \citet{Yang2022} estudian la ventana de 50 a 30 MeV y publican 38.2 MeV para el criterio $n_{\bar B}/(n_B-n_{\bar B})=1$. Estos hitos describen la evolución de las abundancias; 50 MeV no es una medición del primer evento de aniquilación y el segundo criterio no exige que los antibariones hayan desaparecido.

El gas de resonancias hadrónicas, HRG por sus siglas en inglés, permite calcular poblaciones de equilibrio a partir de masas de reposo, espín y cargas de las especies de partícula. Las listas de Thermal-FIST proporcionan un inventario documentado para ese cálculo \citep{Vovchenko2019,TFLists2026}. En este trabajo, las propiedades tabuladas constituyen las firmas utilizadas para identificar cada especie de partícula. No se analizaron señales espectrales originales de la época temprana. La abundancia cosmológica actual se incorporó como restricción externa del número bariónico neto.

El objetivo fue comprobar si esas propiedades, combinadas con la abundancia actual y una temperatura fundamentada en la literatura, permiten inferir por separado las poblaciones tempranas de bariones y antibariones, cuantificar su relación con el residuo y contrastar el resultado con lo publicado. Se evaluó también si el presente determina una reconstrucción histórica única. El análisis se centró en 50 MeV y en la localización del cruce publicado, conservando las masas del inventario HRG y declarando las condiciones que hacen posible la inferencia.
